% Figura 5.3 - schema a blocchi del metodo Newton-CG per problemi L1-regolarizzati
\begin{figure}[h]
\centering
\resizebox{\textwidth}{!}{%
\begin{tikzpicture}[
  node distance=1.1cm and 2.4cm,
  every node/.style={font=\small},
  box/.style={
    draw=coverblue, fill=coverblue!6, rounded corners=3pt,
    text width=4.2cm, align=center, minimum height=1.2cm,
    inner sep=6pt
  },
  io/.style={
    draw=gray!60, fill=gray!8, rounded corners=3pt,
    text width=4.0cm, align=center, minimum height=1.0cm,
    inner sep=5pt
  },
  arrow/.style={-{Stealth[length=2.5mm]}, thick, color=coverblue},
  label/.style={font=\scriptsize, fill=white, inner sep=1.5pt}
]

  % Nodi principali (colonna centrale)
  \node[io] (init) {Inizializzazione\\[2pt] $w_0=0$, $\mathcal{S}_0$, $\mathcal{H}_0$, $\nu$, $\eta$, $\sigma$, $R$};
  \node[box, below=of init, text width=4.6cm, minimum height=1.4cm] (subgrad) {Calcola $\widetilde{\nabla}F_{\mathcal{S}_k}(w_k)$\\[2pt] e determina $z_k$, $\mathcal{A}_k$};
  \node[box, below=of subgrad, text width=4.8cm, minimum height=1.6cm] (cg) {CG nel sottospazio libero\\[2pt] $[Y_k^T H_k Y_k]\,d^Y = -Y_k^T \widetilde{\nabla}F$\\[1pt] Arresto: $\|r\| \le \eta\|Y_k^T \widetilde{\nabla}F\|$};
  \node[box, below=of cg, text width=4.6cm, minimum height=1.4cm] (ls) {Ricerca lineare proiettata\\[2pt] Trova $\alpha_k$\\[1pt] $w_{k+1} = P[w_k + \alpha_k d_k]$};
  \node[box, below=of ls] (recamp) {Re-campionamento\\[2pt] $\mathcal{S}_{k+1}$, $\mathcal{H}_{k+1}$ (fissi)};
  \node[io, below=of recamp] (conv) {Convergenza?\\[2pt] $\|\widetilde{\nabla}F_{\mathcal{S}_k}(w_k)\| < \mathrm{tol}$};
  \node[io, right=2.1cm of conv] (end) {Fine:\\[2pt] $w^* \approx w_k$};

  % Frecce verticali principali
  \draw[arrow] (init) -- (subgrad);
  \draw[arrow] (subgrad) -- (cg);
  \draw[arrow] (cg) -- (ls);
  \draw[arrow] (ls) -- (recamp);
  \draw[arrow] (recamp) -- (conv);

  % Freccia uscita
  \draw[arrow] (conv) -- (end) node[label, midway, above] {sì};

  % Loop ritorno da conv a subgrad (percorso esterno sinistro)
  \draw[arrow] (conv.west) -- ++(-1.6,0) |- (subgrad.west)
    node[label, pos=0.25, left] {no: $k \leftarrow k{+}1$};

\end{tikzpicture}}
\caption{Schema del metodo Newton-CG per problemi con regolarizzazione $L_1$:
il gradiente generalizzato determina l'active set e la faccia ortante, il CG
risolve il sistema ridotto sulle variabili libere e la ricerca lineare
proiettata mantiene le variabili dell'active set a zero.}
\label{fig:newton_l1}
\end{figure}
